\subsection{Superimposed integrals of functions with values in a Banach space}

THEOREM 1. --- Let f be a function with values in a Banach space F or in \(\overline{R}\) , and let H be the set of \(t \in T\) for which f is not \(\lambda_{t}\) -integrable.

\begin{enumerate}
\def\labelenumi{\alph{enumi})}
\tightlist
\item
  If \(\mathbf{f}\) is \(\nu\) -integrable, then \(\mathrm{H}\) is locally \(\mu\) -negligible, the function \(t \mapsto \int \mathbf{f}(x)d\lambda_t(x)\) (defined for \(t \notin \mathrm{H}\) ) is essentially \(\mu\) -integrable, and
\end{enumerate}

\[
\int \mathbf {f} (x) d \nu (x) = \int d \mu (t) \int \mathbf {f} (x) d \lambda_ {t} (x).\tag{11}
\]

\begin{enumerate}
\def\labelenumi{\alph{enumi})}
\setcounter{enumi}{1}
\item
  If \(\mathbf{f}\) is \(\nu\) -integrable and \(\Lambda\) is vaguely continuous, then \(\mathbf{H}\) is moreover \(\mu\) -negligible and the function \(t \mapsto \int \mathbf{f}(x) d\lambda_t(x)\) (defined for \(t \notin \mathbf{H}\) ) is \(\mu\) -integrable.
\item
  If \(\lambda_t^\bullet(1) < +\infty\) locally \(\mu\) -almost everywhere, then the conclusions of a) remain true for f an essentially \(\nu\) -integrable function.
\end{enumerate}

We are first going to establish a) (resp. b)). This statement is true when f is a positive numerical function (Prop. 5); if f is an integrable function with values in \(\overline{R}\) , this result may be applied to the positive functions \(f^{+}\) and \(f^{-}\) , and therefore extends at once to f by subtraction. It remains to treat the case of functions with values in F. Let H be the subspace of \(\mathcal{L}_{\mathrm{F}}^{1}(\nu)\) formed by the linear combinations, with coefficients in F, of the functions in \(\mathcal{H}(\mathrm{X})\) ; the result pertaining to real functions implies at once the validity of the statement for the elements of H. Now, H is dense in \(\mathcal{L}_{\mathrm{F}}^{1}(\nu)\) ; therefore, for every \(\mathbf{f} \in \mathcal{L}_{\mathrm{F}}^{1}(\nu)\) , there exists a sequence \((\mathbf{f}_{n})\) of elements of H that has the following properties:

\begin{enumerate}
\def\labelenumi{\arabic{enumi})}
\item
  the sequence \((\mathbf{f}_n)\) converges to \(\mathbf{f}\) in mean in \(\mathcal{L}_{\mathrm{F}}^{1}(\nu)\) , and \(\nu\) -almost everywhere;
\item
  the function \(g = |\mathbf{f}_0| + \sum_{n\in \mathbf{N}}|\mathbf{f}_{n + 1} - \mathbf{f}_n|\) is such that \(\nu^{*}(g) < +\infty\) (Ch. IV, §3, No.~4, Th. 3).
\end{enumerate}

Let \(N_{1}\) be the set of \(t \in T\) such that \(\lambda_{t}^{*}(g) = +\infty\) ; \(N_{1}\) is locally \(\mu\) -negligible (resp. \(\mu\) -negligible) by formula (6) (resp. (7)). For \(t \notin N_{1}\) , the \(f_{n}\) belong to \(\mathcal{L}_{F}^{1}(\lambda_{t})\) , the sequence \((f_{n})\) converges \(\lambda_{t}\) -almost everywhere, as well as for the topology of convergence in mean in \(\mathcal{L}_{F}^{1}(\lambda_{t})\) (Ch. IV, §3, No.~3, Prop. 6). Let M be set of \(x \in X\) such that \(f_{n}(x)\) does not converge to \(f(x)\) : since M is \(\nu\) -negligible, the set \(N_{2}\) of \(t \in T\) such that M is not \(\lambda_{t}\) -negligible is locally \(\mu\) -negligible (resp. \(\mu\) -negligible) by Cor. 1 of Prop. 3.

Suppose that \(t\) does not belong to \(\mathrm{N}_1 \cup \mathrm{N}_2\) ; the sequence \((\mathbf{f}_n)\) converges in mean in \(\mathcal{L}_{\mathrm{F}}^{1}(\lambda_t)\) , and converges \(\lambda_t\) -almost everywhere to \(\mathbf{f}\) . Therefore \(\mathbf{f} \in \mathcal{L}_{\mathrm{F}}^{1}(\lambda_t)\) and \(\int \mathbf{f} d\lambda_t = \lim_{n \to \infty} \int \mathbf{f}_n d\lambda_t\) (Ch. IV, §4, No.~1). The set H of the statement is thus contained in \(N_{1} \cup N_{2}\) ; it is therefore locally \(\mu\) -negligible (resp. \(\mu\) -negligible). On the other hand, the function \(t \mapsto \int f d\lambda_{t}\) is equal locally \(\mu\) -almost everywhere to the limit of a sequence of \(\mu\) -measurable functions; it is therefore \(\mu\) -measurable. Finally, for every \(t \notin N_{1} \cup N_{2}\) and every n, we have

\[
\left| \int \mathbf {f} _ {n} (x) d \lambda_ {t} (x) \right| \leqslant \int^ {*} g (x) d \lambda_ {t} (x)
\]

by virtue of the inequality \(|\mathbf{f}_n| \leqslant g\) and Prop. 2 of Ch. IV, §4, No.~2. Now, the function \(t \mapsto \int^{*} g(x) d\lambda_t(x)\) is essentially \(\mu\) -integrable (resp. \(\mu\) -integrable) by Prop. 5. We may therefore apply Lebesgue's theorem, which yields

\[
\int d \mu (t) \int \mathbf {f} (x) d \lambda_ {t} (x) = \lim _ {n \rightarrow \infty} \int d \mu (t) \int \mathbf {f} _ {n} (x) d \lambda_ {t} (x) = \lim _ {n \rightarrow \infty} \int \mathbf {f} _ {n} (x) d \nu (x).
\]

Since \(\int \mathbf{f}_n(x) d\nu(x)\) tends to \(\int \mathbf{f}(x) d\nu(x)\) as \(n\) tends to \(\infty\) , by the hypotheses made on the sequence \((\mathbf{f}_n)\) , the relation (11) follows and we have proved a) (resp. b)).

Now suppose that \(\lambda_t^\bullet(1) < +\infty\) locally \(\mu\) -almost everywhere, and that \(\mathbf{g}\) is an essentially \(\nu\) -integrable function. Let \(\mathbf{f}\) be a \(\nu\) -integrable function such that \(\mathbf{g} = \mathbf{f}\) locally \(\nu\) -almost everywhere (§1, No.~3). Then \(\mathbf{g} = \mathbf{f}\) almost everywhere for \(\lambda_t\) , except for \(t\) forming a locally \(\mu\) -negligible set \(\mathrm{P}\) (Cor. 1 c) of Prop. 3). Therefore \(\int \mathbf{g} d\lambda_t = \int \mathbf{f} d\lambda_t\) for all \(t \notin \mathrm{P} \cup \mathrm{H}\) , and this completes the proof.

Remark. --- Let \(\Lambda: t \mapsto \lambda_{t}\) be a \(\mu\) -adequate mapping of T into \(\mathcal{M}_{+}(X)\) . If a mapping \(\Lambda': t \mapsto \lambda_{t}'\) of T into \(\mathcal{M}_{+}(X)\) is equal to \(\Lambda\) locally \(\mu\) -almost everywhere, it follows at once from the definitions that \(\Lambda'\) is also \(\mu\) -adequate, and that \(\Lambda\) and \(\Lambda'\) have the same integral. If now \(H: t \mapsto \eta_{t}\) is a function with values in \(\mathcal{M}_{+}(X)\) , defined locally \(\mu\) -almost everywhere, we shall again say that H is \(\mu\) -adequate if it is equal locally \(\mu\) -almost everywhere to a mapping \(\Lambda: t \mapsto \lambda_{t}\) that is everywhere defined and \(\mu\) -adequate. We then set \(\int \eta_{t} d\mu(t) = \int \lambda_{t} d\mu(t)\) , a definition that does not depend on the function \(\Lambda\) utilized. We leave to the reader the task of verifying that the propositions proved in the preceding Nos. extend to \(\mu\) -adequate functions defined locally \(\mu\) -almost everywhere.

